% ============================================================================
% LANGENTWURF-TEMPLATE - Greenhouse Schools Rostock
% ============================================================================
% Konform mit: Fächerübergreifende Handreichung M-V
% Version: 2.0 (Dezember 2024)
%
% TIER-SYSTEM:
%   Tier 1 (FULL)     = Alle Abschnitte (Prüfungsstunden, komplexe Themen)
%   Tier 2 (STANDARD) = Ohne Reflexion, verkürzte Sachanalyse
%   Tier 3 (MINIMAL)  = Nur Ziele, Verlaufsplan, Lösungen
%
% LERNGRUPPEN-PROFILE:
%   A = Gesamtschule heterogen (BR+MR+GY, Inklusion) - Kl. 7-9
%   B = E-Kurs Gesamtschule (GY-Tendenz) - Kl. 9-10
%   C = G-Kurs Gesamtschule (MR-Tendenz) - Kl. 9-10
%   D = Gymnasium Sek I - Kl. 7-10
%   E = Gymnasium Sek II / Oberstufe - Kl. 11-12
%
% TESTDIFFERENZIERUNG:
%   Physik, Chemie: 3 Niveaus (BR/MR/GY oder */**/***)
%   Informatik, Spanisch: 1 Niveau (einheitlich)
%
% VERWENDUNG:
%   1. Kopiere diese Datei als LE-XX-thema.tex
%   2. Setze \plantier{1}, {2} oder {3}
%   3. Setze \lerngruppenprofil{A}, {B}, {C}, {D} oder {E}
%   4. Fülle die Variablen aus
%   5. Kompiliere 2x mit pdflatex
% ============================================================================

\documentclass[11pt,a4paper]{article}

% ============================================================================
% PACKAGES
% ============================================================================
\usepackage[utf8]{inputenc}
\usepackage[T1]{fontenc}
\usepackage[ngerman]{babel}
\usepackage{geometry}
\usepackage{fancyhdr}
\usepackage{lastpage}
\usepackage{xcolor}
\usepackage{tcolorbox}
\usepackage{enumitem}
\usepackage{tabularx}
\usepackage{longtable}
\usepackage{array}
\usepackage{booktabs}
\usepackage{hyperref}
\usepackage{ifthen}
\usepackage{tikz}
\usepackage{amssymb}

% ============================================================================
% KONFIGURATION
% ============================================================================

% --- TIER (1, 2 oder 3) ---
\newcommand{\plantier}{1}

% --- LERNGRUPPEN-PROFIL (A, B, C, D oder E) ---
\newcommand{\lerngruppenprofil}{A}

% --- FACH (physik, chemie, informatik, spanisch) ---
\newcommand{\fachid}{spanisch}

% Bedingte Befehle basierend auf Tier
\newcommand{\tierone}[1]{\ifthenelse{\equal{\plantier}{1}}{#1}{}}
\newcommand{\tieronetwo}[1]{\ifthenelse{\equal{\plantier}{1}\OR\equal{\plantier}{2}}{#1}{}}
\newcommand{\tierall}[1]{#1}

% Bedingte Befehle basierend auf Fach (Test-Differenzierung)
\newcommand{\testdifferenziert}[1]{%
    \ifthenelse{\equal{\fachid}{physik}\OR\equal{\fachid}{chemie}}{#1}{}%
}
\newcommand{\testeinheitlich}[1]{%
    \ifthenelse{\equal{\fachid}{informatik}\OR\equal{\fachid}{spanisch}}{#1}{}%
}

% Bedingte Befehle basierend auf Profil (Heterogenität)
\newcommand{\profilheterogen}[1]{%
    \ifthenelse{\equal{\lerngruppenprofil}{A}}{#1}{}%
}
\newcommand{\profilhomogen}[1]{%
    \ifthenelse{\equal{\lerngruppenprofil}{D}\OR\equal{\lerngruppenprofil}{E}}{#1}{}%
}

% ============================================================================
% VARIABLEN - HIER AUSFÜLLEN
% ============================================================================

% --- Metadaten ---
\newcommand{\fach}{Spanisch}                    % Anzeigename
\newcommand{\jahrgangsstufe}{11}                % 5-12
\newcommand{\schulart}{Gymnasium}               % Gymnasium, Gesamtschule
\newcommand{\kursart}{Grundkurs}                % Grundkurs, Leistungskurs, E-Kurs, G-Kurs
\newcommand{\stundenthema}{Das Verb gustar}     % Konkretes Thema der Stunde
\newcommand{\reihenthema}{Mis gustos}           % Übergeordnetes Reihenthema
\newcommand{\stundennummer}{1}                  % X. Stunde der Reihe
\newcommand{\reihenstunden}{6}                  % Gesamtzahl Stunden in Reihe
\newcommand{\unterrichtsdatum}{XX.XX.XXXX}      % Datum der Durchführung
\newcommand{\unterrichtszeit}{XX:XX -- XX:XX}   % Zeitrahmen
\newcommand{\raum}{XXX}                         % Raumnummer

% --- Lehrkraft ---
\newcommand{\lehrkraft}{N.N.}                   % Name der Lehrkraft
\newcommand{\ausbildungsstand}{LAA}             % LAA, Studienreferendar, Lehrkraft

% --- Lerngruppe (wird durch Profil vordefiniert, kann überschrieben werden) ---
\newcommand{\lerngruppengroesse}{24}            % Anzahl SuS
\newcommand{\maennlich}{12}                     % Anzahl männlich
\newcommand{\weiblich}{12}                      % Anzahl weiblich
\newcommand{\divers}{0}                         % Anzahl divers

% --- Kompetenzen/Niveau ---
\newcommand{\gerniveau}{A2}                     % A1, A2, B1, B2 (für Sprachen)

% --- Fachfarbe (Hex ohne #) ---
\newcommand{\fachfarbe}{c41e3a}
% Physik: 2c5aa0, Chemie: 2a7a4b, Informatik: b35c00, Spanisch: c41e3a

% ============================================================================
% LERNGRUPPEN-PROFILE (FIKTIVE REFERENZKLASSEN)
% ============================================================================
% Diese Profile definieren typische Klassenzusammensetzungen für die Planung.
% Sie decken alle realistischen Szenarien ab und ermöglichen konsistente
% Materialentwicklung.

\newcommand{\profilbeschreibung}{%
\ifthenelse{\equal{\lerngruppenprofil}{A}}{%
    % PROFIL A: Gesamtschule heterogen (Kl. 7-9)
    \textbf{Profil A: Gesamtschule heterogen}\\[0.5em]
    \begin{tabular}{ll}
        Klassenstufe: & 7--9 \\
        Zusammensetzung: & BR (20\%), MR (50\%), GY (30\%) \\
        Inklusion: & 2--3 SuS mit LRS, 1 SuS mit Förderbedarf \\
        Testdifferenzierung: & 3 Niveaus (*/**/***) \\
        Besonderheiten: & Hohe Heterogenität, intensive Differenzierung nötig
    \end{tabular}
}{}%
\ifthenelse{\equal{\lerngruppenprofil}{B}}{%
    % PROFIL B: E-Kurs Gesamtschule (Kl. 9-10)
    \textbf{Profil B: E-Kurs Gesamtschule}\\[0.5em]
    \begin{tabular}{ll}
        Klassenstufe: & 9--10 \\
        Zusammensetzung: & MR (30\%), GY (70\%) \\
        Inklusion: & Ggf. 1 SuS mit LRS (Nachteilsausgleich) \\
        Testdifferenzierung: & 1--2 Niveaus (Standard + Zusatz) \\
        Besonderheiten: & Leistungsorientiert, GY-Vorbereitung
    \end{tabular}
}{}%
\ifthenelse{\equal{\lerngruppenprofil}{C}}{%
    % PROFIL C: G-Kurs Gesamtschule (Kl. 9-10)
    \textbf{Profil C: G-Kurs Gesamtschule}\\[0.5em]
    \begin{tabular}{ll}
        Klassenstufe: & 9--10 \\
        Zusammensetzung: & BR (40\%), MR (60\%) \\
        Inklusion: & 2--4 SuS mit LRS, teils Förderbedarf \\
        Testdifferenzierung: & 2 Niveaus (Basis + Standard) \\
        Besonderheiten: & Praxisorientiert, MR-Vorbereitung
    \end{tabular}
}{}%
\ifthenelse{\equal{\lerngruppenprofil}{D}}{%
    % PROFIL D: Gymnasium Sek I (Kl. 7-10)
    \textbf{Profil D: Gymnasium Sek I}\\[0.5em]
    \begin{tabular}{ll}
        Klassenstufe: & 7--10 \\
        Zusammensetzung: & 100\% GY-Niveau \\
        Inklusion: & Ggf. einzelne SuS mit LRS (Nachteilsausgleich) \\
        Testdifferenzierung: & 1 Niveau (mit BE-Differenzierung) \\
        Besonderheiten: & Homogen, Abitur-Vorbereitung
    \end{tabular}
}{}%
\ifthenelse{\equal{\lerngruppenprofil}{E}}{%
    % PROFIL E: Gymnasium Sek II / Oberstufe (Kl. 11-12)
    \textbf{Profil E: Gymnasium Sek II / Oberstufe}\\[0.5em]
    \begin{tabular}{ll}
        Klassenstufe: & 11--12 \\
        Zusammensetzung: & 100\% Abiturniveau \\
        Kursart: & Grundkurs / Leistungskurs \\
        Testdifferenzierung: & 1 Niveau (15-NP-Skala) \\
        Besonderheiten: & Wissenschaftspropädeutisch, EPA-konform
    \end{tabular}
}{}%
}

% ============================================================================
% LAYOUT
% ============================================================================
\geometry{left=2.5cm, right=2.5cm, top=2.5cm, bottom=2.5cm}
\definecolor{fach}{HTML}{\fachfarbe}
\definecolor{gray}{HTML}{666666}
\definecolor{lightgray}{HTML}{f5f5f5}
\definecolor{successgreen}{HTML}{2a7a4b}
\definecolor{warningorange}{HTML}{b35c00}

\tcbuselibrary{skins,breakable}

% Farbige Boxen
\newtcolorbox{infobox}[1][]{
    colback=lightgray,
    colframe=fach,
    fonttitle=\bfseries,
    title=#1,
    breakable
}

\newtcolorbox{scorebox}{
    colback=white,
    colframe=fach,
    fonttitle=\bfseries,
    title=Didaktik-Score,
    breakable
}

\newtcolorbox{profilbox}{
    colback=lightgray,
    colframe=gray,
    fonttitle=\bfseries,
    title=Lerngruppen-Profil,
    breakable
}

% Kopf-/Fußzeile
\pagestyle{fancy}
\fancyhf{}
\fancyhead[L]{\small\textcolor{gray}{\fach{} -- Jg. \jahrgangsstufe{} -- \stundenthema}}
\fancyhead[R]{\small\textcolor{gray}{Langentwurf Tier \plantier{} / Profil \lerngruppenprofil}}
\fancyfoot[C]{\small\thepage{} / \pageref{LastPage}}
\renewcommand{\headrulewidth}{0.4pt}
\renewcommand{\footrulewidth}{0pt}

% Abschnitte
\renewcommand{\thesection}{\arabic{section}}
\renewcommand{\thesubsection}{\thesection.\arabic{subsection}}

% ============================================================================
% DOKUMENT
% ============================================================================
\begin{document}

% ============================================================================
% DECKBLATT (Tier 1+2)
% ============================================================================
\tieronetwo{
\begin{titlepage}
    \centering
    \vspace*{2cm}

    {\Large\textcolor{gray}{Greenhouse Schools Rostock}}\\[0.5cm]
    {\large\textcolor{gray}{\schulart}}\\[2cm]

    {\Huge\textcolor{fach}{\textbf{Langentwurf}}}\\[0.5cm]
    {\large Tier \plantier{} --
        \ifthenelse{\equal{\plantier}{1}}{Vollständig}{%
        \ifthenelse{\equal{\plantier}{2}}{Standard}{Minimal}}
        \quad|\quad Profil \lerngruppenprofil}\\[2cm]

    {\LARGE\textbf{\stundenthema}}\\[0.5cm]
    {\large im Rahmen der Unterrichtsreihe}\\[0.3cm]
    {\Large\textit{\reihenthema}}\\[2cm]

    \begin{tabular}{rl}
        \textbf{Fach:} & \fach \\[0.3cm]
        \textbf{Klasse/Kurs:} & \jahrgangsstufe{} (\kursart) \\[0.3cm]
        \textbf{Datum:} & \underline{\hspace{3cm}} \\[0.3cm]
        \textbf{Zeit:} & \underline{\hspace{3cm}} \\[0.3cm]
        \textbf{Raum:} & \underline{\hspace{3cm}} \\[0.3cm]
    \end{tabular}

    \vfill

    {\large Lehrkraft: \lehrkraft}\\[0.3cm]
    {\normalsize\textcolor{gray}{\ausbildungsstand}}\\[1cm]

\end{titlepage}
}

% ============================================================================
% INHALTSVERZEICHNIS (Tier 1+2)
% ============================================================================
\tieronetwo{
\tableofcontents
\newpage
}

% ============================================================================
% KONFIGURATIONS-ÜBERSICHT
% ============================================================================
\begin{infobox}[Konfiguration dieses Langentwurfs]
\begin{tabular}{ll}
    \textbf{Tier:} & \plantier{} --
        \ifthenelse{\equal{\plantier}{1}}{Vollständig (alle Abschnitte)}{%
        \ifthenelse{\equal{\plantier}{2}}{Standard (ohne Reflexion)}{Minimal (Ziele + Verlauf)}} \\[0.3em]
    \textbf{Profil:} & \lerngruppenprofil{} --
        \ifthenelse{\equal{\lerngruppenprofil}{A}}{Gesamtschule heterogen}{%
        \ifthenelse{\equal{\lerngruppenprofil}{B}}{E-Kurs Gesamtschule}{%
        \ifthenelse{\equal{\lerngruppenprofil}{C}}{G-Kurs Gesamtschule}{%
        \ifthenelse{\equal{\lerngruppenprofil}{D}}{Gymnasium Sek I}{Gymnasium Sek II}}}} \\[0.3em]
    \textbf{Fach:} & \fach{} (\fachid) \\[0.3em]
    \textbf{Testdifferenzierung:} &
        \ifthenelse{\equal{\fachid}{physik}\OR\equal{\fachid}{chemie}}{%
            Ja (3 Niveaus: */**/***)}{%
            Nein (einheitliches Niveau)}
\end{tabular}
\end{infobox}

% ============================================================================
% LERNGRUPPEN-PROFIL (Tier 1+2)
% ============================================================================
\tieronetwo{
\begin{profilbox}
\profilbeschreibung
\end{profilbox}
}

% ============================================================================
% 1. BEDINGUNGSANALYSE (Tier 1+2)
% ============================================================================
\tieronetwo{
\section{Bedingungsanalyse}

\subsection{Statistik der Lerngruppe}
\begin{tabular}{ll}
    Kursgröße: & \lerngruppengroesse{} SuS (\maennlich{} m / \weiblich{} w / \divers{} d) \\
    Jahrgangsstufe: & \jahrgangsstufe \\
    Kursart: & \kursart \\
    Profil: & \lerngruppenprofil \\
\end{tabular}

\subsection{Zusammensetzung nach Leistungsniveau}
\profilheterogen{%
\begin{tabular}{|l|c|l|}
\hline
\textbf{Niveau} & \textbf{Anteil} & \textbf{Charakteristik} \\
\hline
Berufsreife (BR / *) & ca. 20\% & Qualitative Beschreibungen, Kopfrechnen \\
\hline
Mittlere Reife (MR / **) & ca. 50\% & Formelanwendung, mehrstufige Aufgaben \\
\hline
Gymnasium (GY / ***) & ca. 30\% & Transfer, Analyse, komplexe Probleme \\
\hline
\end{tabular}
}
\profilhomogen{%
\begin{tabular}{|l|c|l|}
\hline
\textbf{Niveau} & \textbf{Anteil} & \textbf{Charakteristik} \\
\hline
Gymnasium (GY) & 100\% & Einheitliches Anforderungsniveau \\
\hline
\end{tabular}
}

\subsection{Individuelle Besonderheiten (Inklusion)}
% TODO: Ausfüllen (anonymisiert)
\profilheterogen{%
\begin{itemize}
    \item \textbf{LRS:} 2--3 SuS (Nachteilsausgleich: Zeitzugabe, vergrößerte Schrift)
    \item \textbf{Förderbedarf Lernen:} 1 SuS (reduzierte Aufgabenmenge, Hilfsmittel)
    \item \textbf{Besonders leistungsstark:} 2--3 SuS (Zusatzaufgaben, Expertenrolle)
    \item \textbf{DaZ:} ggf. 1--2 SuS (Wortschatzarbeit, Sprachhilfen)
\end{itemize}
}
\profilhomogen{%
\begin{itemize}
    \item \textbf{LRS:} Ggf. Einzelfälle mit Nachteilsausgleich
    \item \textbf{Besonders leistungsstark:} Differenzierung nach oben durch Zusatzaufgaben
\end{itemize}
}

\subsection{Arbeitsvoraussetzungen}
\begin{itemize}
    \item Raumsituation: ...
    \item Technische Ausstattung: Beamer, WLAN, Tablets/PCs (optional)
    \item Verfügbare Medien: Tafel, Arbeitsblätter, digitaler Lernpfad
\end{itemize}

\subsection{Fachbezogener Entwicklungsstand}
\begin{itemize}
    \item Vorwissen: ...
    \item Bekannte Methoden: Einzelarbeit, Partnerarbeit, Plenum
    \ifthenelse{\equal{\fachid}{spanisch}\OR\equal{\fachid}{franzoesisch}}{%
        \item Sprachniveau: \gerniveau
    }{}
\end{itemize}

\subsection{Erfahrungen mit der Lerngruppe}
\begin{itemize}
    \item Bisherige Unterrichtserfahrungen: ...
    \item Bewährte Sozialformen: ...
    \item Herausforderungen: ...
\end{itemize}
}

% ============================================================================
% 2. SACHANALYSE (Tier 1 vollständig, Tier 2 verkürzt)
% ============================================================================
\tieronetwo{
\section{Sachanalyse}

\tierone{
\subsection{Fachwissenschaftliche Einordnung}
% TODO: Ausführliche fachliche Darstellung
\begin{itemize}
    \item Systematische Einordnung des Themas
    \item Fachliche Grundlagen
    \item Aktuelle Forschung/Entwicklungen (wo relevant)
\end{itemize}

\subsection{Kontrastive Betrachtung}
% Bei Sprachen: Vergleich mit Muttersprache (Interferenzen)
% Bei MINT: Alltagsvorstellungen vs. Fachkonzepte
\begin{itemize}
    \item Typische Fehlvorstellungen / Interferenzen
    \item Transfermöglichkeiten (positiv/negativ)
\end{itemize}
}

\subsection{Kernkonzepte der Stunde}
% TODO: Die 3-5 wichtigsten fachlichen Punkte
\begin{enumerate}
    \item Konzept 1: ...
    \item Konzept 2: ...
    \item Konzept 3: ...
\end{enumerate}

\tierone{
\subsection{Weiterführende Aspekte}
% Was kommt in späteren Stunden? Vertikale Vernetzung
\begin{itemize}
    \item Anknüpfungspunkte für Folgestunden: ...
    \item Vertiefungsmöglichkeiten: ...
\end{itemize}
}
}

% ============================================================================
% 3. DIDAKTISCHE ANALYSE (Tier 1+2)
% ============================================================================
\tieronetwo{
\section{Didaktische Analyse}

\subsection{Stellung der Stunde in der Unterrichtseinheit}
Dies ist die \textbf{\stundennummer. Stunde} der Unterrichtsreihe ``\reihenthema'' (\reihenstunden{} Stunden gesamt).

\begin{tabular}{|c|l|l|}
\hline
\textbf{Std.} & \textbf{Thema} & \textbf{Schwerpunkt} \\
\hline
1 & ... & Einführung \\
2 & ... & Vertiefung \\
... & ... & ... \\
\hline
\end{tabular}

\subsection{Ziele der Unterrichtseinheit}
Am Ende der Reihe können die SuS...
\begin{itemize}
    \item ...
    \item ...
    \item ...
\end{itemize}

\subsection{Legitimation}
\textbf{Rahmenplanbezug (M-V):}
\begin{itemize}
    \item Kompetenzbereich: ...
    \item Standard/Kompetenzerwartung: ...
    \item Inhaltsfeld: ...
\end{itemize}

\textbf{Bildungswert (Klafki):}
\begin{itemize}
    \item Gegenwartsbedeutung: ...
    \item Zukunftsbedeutung: ...
    \item Exemplarische Bedeutung: ...
\end{itemize}

\subsection{Didaktische Reduktion}
\textbf{Bewusst ausgeklammert:}
\begin{itemize}
    \item ...
    \item ...
\end{itemize}

\textbf{Begründung:} ...

\subsection{Strukturierung des Unterrichts}
% Phasierung mit Begründung
\begin{enumerate}
    \item \textbf{Einstieg:} ... (Aktivierung, Problemstellung)
    \item \textbf{Erarbeitung:} ... (Input, Entdeckung)
    \item \textbf{Übung/Anwendung:} ... (Festigung, Transfer)
    \item \textbf{Sicherung:} ... (Zusammenfassung, Reflexion)
\end{enumerate}
}

% ============================================================================
% 4. STUNDENZIELE (ALLE TIERS)
% ============================================================================
\section{Stundenziele}

\subsection{Grobziele}
Die SuS erweitern ihre \textbf{[Kompetenz]} durch \textbf{[Handlung]},
indem sie \textbf{[konkrete Aktivität]}.

\tieronetwo{
\textbf{Kompetenzbereiche:}
\begin{itemize}
    \item Fachkompetenz: ...
    \item Methodenkompetenz: ...
    \item Sozialkompetenz: ...
    \item Selbstkompetenz: ...
\end{itemize}
}

\subsection{Feinziele (operationalisiert)}
\begin{tabular}{|c|p{9cm}|c|c|}
\hline
\textbf{Nr.} & \textbf{Die SuS können...} & \textbf{AFB} & \textbf{Niveau} \\
\hline
FZ1 & ... & I & alle \\
\hline
FZ2 & ... & II & alle \\
\hline
FZ3 & ... & II & **/*** \\
\hline
FZ4 & ... & III & *** \\
\hline
\end{tabular}

\vspace{1em}
\textbf{AFB-Verteilung:}
\begin{itemize}
    \item AFB I (Reproduktion): ca. 25\%
    \item AFB II (Reorganisation/Transfer): ca. 50\%
    \item AFB III (Reflexion/Problemlösung): ca. 25\%
\end{itemize}

% ============================================================================
% 5. METHODISCHE ANALYSE (Tier 1+2)
% ============================================================================
\tieronetwo{
\section{Methodische Analyse}

\subsection{Einstieg}
\begin{itemize}
    \item \textbf{Methode:} ...
    \item \textbf{Begründung:} ...
    \item \textbf{Alternative:} ...
\end{itemize}

\subsection{Erarbeitung}
\begin{itemize}
    \item \textbf{Methode:} ...
    \item \textbf{Begründung:} ...
    \item \textbf{Alternative:} ...
\end{itemize}

\subsection{Übungsphasen}
\begin{itemize}
    \item \textbf{Methode:} ...
    \item \textbf{Progression:} von gelenkt zu frei
    \item \textbf{Begründung:} ...
\end{itemize}

\subsection{Sozialformen}
\begin{tabular}{|l|l|l|}
\hline
\textbf{Phase} & \textbf{Sozialform} & \textbf{Begründung} \\
\hline
Einstieg & Plenum & Gemeinsamer Fokus, Aktivierung \\
Erarbeitung & Einzelarbeit & Individuelle Verarbeitung \\
Übung & Partnerarbeit & Kooperatives Lernen, Sprachproduktion \\
Sicherung & Plenum & Ergebnissicherung, Klärung \\
\hline
\end{tabular}

\subsection{Medien}
\begin{itemize}
    \item Tafel/Whiteboard: Tafelbild, Visualisierung
    \item Digitale Medien: Lernpfad (optional), Beamer
    \item Arbeitsblätter: AB-XX-thema.pdf
\end{itemize}

\subsection{Differenzierung}
\testdifferenziert{%
% Für Physik/Chemie: 3 Niveaus
\begin{tabular}{|l|p{3.5cm}|p{3.5cm}|p{3.5cm}|}
\hline
& \textbf{Basis (*)} & \textbf{Standard (**)} & \textbf{Erweitert (***)} \\
\hline
\textbf{Inhalt} & Qualitativ, reduziert & Formelanwendung & Transfer, Analyse \\
\hline
\textbf{Hilfen} & Lückentext, Beispiele & Teillösung, Tipps & Nur Impulse \\
\hline
\textbf{Produkt} & Einfache Antworten & Begründete Aussagen & Eigene Erklärungen \\
\hline
\textbf{Test} & Separates Niveau & Separates Niveau & Separates Niveau \\
\hline
\end{tabular}
}
\testeinheitlich{%
% Für Informatik/Spanisch: 1 Niveau mit interner Differenzierung
\begin{tabular}{|l|p{4.5cm}|p{4.5cm}|}
\hline
& \textbf{Basis} & \textbf{Erweitert} \\
\hline
\textbf{Inhalt} & Kernaufgaben (Pflicht) & Zusatzaufgaben (Kür) \\
\hline
\textbf{Hilfen} & Wortschatz, Beispiele & Nur Impulse \\
\hline
\textbf{Produkt} & Standardanforderung & Kreative Erweiterung \\
\hline
\textbf{Test} & \multicolumn{2}{c|}{Einheitliches Niveau (Punktedifferenzierung)} \\
\hline
\end{tabular}
}

\subsection{Erwartungshorizonte}
\begin{itemize}
    \item \textbf{Minimum:} ... (alle SuS erreichen)
    \item \textbf{Regelstandard:} ... (Mehrheit erreicht)
    \item \textbf{Optimal:} ... (leistungsstarke SuS)
\end{itemize}

\subsection{Überprüfung der Teilziele}
\begin{tabular}{|c|l|l|}
\hline
\textbf{FZ} & \textbf{Überprüfung durch} & \textbf{Zeitpunkt} \\
\hline
FZ1 & ... & ... \\
FZ2 & ... & ... \\
FZ3 & ... & ... \\
FZ4 & ... & ... \\
\hline
\end{tabular}
}

% ============================================================================
% 6. VERLAUFSPLANUNG (ALLE TIERS)
% ============================================================================
\section{Verlaufsplanung}

\begin{longtable}{|p{1cm}|p{1.5cm}|p{4cm}|p{3cm}|p{2cm}|p{2cm}|}
\hline
\textbf{Zeit} & \textbf{Phase} & \textbf{Inhalt / Lehrerhandeln} & \textbf{Schülerhandeln} & \textbf{Sozial-form} & \textbf{Medien} \\
\hline
\endhead

0--5 & Einstieg & ... & ... & Plenum & ... \\
\hline

5--15 & Erarbei-tung & ... & ... & EA & ... \\
\hline

15--25 & Übung 1 & ... & ... & PA & ... \\
\hline

25--35 & Übung 2 & ... & ... & GA & ... \\
\hline

35--40 & Siche-rung & ... & ... & Plenum & ... \\
\hline

40--45 & Abschluss & ... & ... & Plenum & ... \\
\hline
\end{longtable}

\textbf{Legende:} EA = Einzelarbeit, PA = Partnerarbeit, GA = Gruppenarbeit

% ============================================================================
% 7. REFLEXION (NUR Tier 1, nach Durchführung)
% ============================================================================
\tierone{
\section{Reflexion (nach der Durchführung)}

\textit{Dieser Abschnitt wird nach der Unterrichtsdurchführung ausgefüllt.}

\subsection{Realisierung der Lernziele}
\begin{tabular}{|c|c|l|}
\hline
\textbf{Feinziel} & \textbf{Erreicht?} & \textbf{Evidenz/Anmerkung} \\
\hline
FZ1 & $\Box$ ja / $\Box$ teilweise / $\Box$ nein & ... \\
\hline
FZ2 & $\Box$ ja / $\Box$ teilweise / $\Box$ nein & ... \\
\hline
FZ3 & $\Box$ ja / $\Box$ teilweise / $\Box$ nein & ... \\
\hline
FZ4 & $\Box$ ja / $\Box$ teilweise / $\Box$ nein & ... \\
\hline
\end{tabular}

\subsection{Wirksamkeit des Vorgehens}
\begin{itemize}
    \item Was hat gut funktioniert? ...
    \item Was hat nicht funktioniert? ...
    \item Unvorhergesehene Ereignisse: ...
\end{itemize}

\subsection{Konsequenzen für die Weiterführung}
\begin{itemize}
    \item Für die nächste Stunde: ...
    \item Für ähnliche Stunden: ...
    \item Für die Reihenplanung: ...
\end{itemize}

\subsection{Gewonnene Einsichten}
\begin{itemize}
    \item Fachdidaktisch: ...
    \item Methodisch: ...
    \item Pädagogisch: ...
\end{itemize}
}

% ============================================================================
% 8. ANLAGEN (ALLE TIERS)
% ============================================================================
\section{Anlagen}

\begin{enumerate}
    \item Arbeitsblatt (AB-XX-thema.pdf)
    \item Musterlösung (ML-XX-thema.pdf)
    \item Lernpfad (LP-XX-thema.html) -- optional
    \item Lehrerhinweise (LH-XX-thema.pdf)
    \item Einstiegsfolie / Tafelbild
    \item ggf. weitere Materialien
\end{enumerate}

% ============================================================================
% 9. LÖSUNGEN (ALLE TIERS)
% ============================================================================
\section{Lösungen}

\subsection{Lösungen Arbeitsblatt}
\begin{enumerate}
    \item Aufgabe 1: ...
    \item Aufgabe 2: ...
    \item Aufgabe 3: ...
\end{enumerate}

\subsection{Erwartete Schülerantworten}
\begin{itemize}
    \item Frage X: Mögliche Antworten sind...
\end{itemize}

% ============================================================================
% 10. LITERATUR (Tier 1+2)
% ============================================================================
\tieronetwo{
\section{Literatur}

\subsection{Fachliteratur}
\begin{itemize}
    \item Autor (Jahr): Titel. Verlag.
\end{itemize}

\subsection{Didaktische Literatur}
\begin{itemize}
    \item Autor (Jahr): Titel. Verlag.
\end{itemize}

\subsection{Rahmenplan}
\begin{itemize}
    \item Ministerium für Bildung, Wissenschaft und Kultur M-V (Jahr): Rahmenplan \fach. Schwerin.
\end{itemize}

\subsection{Digitale Quellen}
\begin{itemize}
    \item URL (Zugriff: Datum)
\end{itemize}
}

% ============================================================================
% 11. DIDAKTIK-SCORE (ALLE TIERS)
% ============================================================================
\newpage
\section{Didaktik-Score}

\begin{scorebox}
\textbf{Bewertung der didaktischen Qualität nach 10 Dimensionen}\\
\textbf{Ziel-Score: $\geq$ 9.0 (Premium-Qualität)}

\vspace{1em}
\begin{tabular}{|c|l|c|c|p{4.5cm}|}
\hline
\textbf{\#} & \textbf{Dimension} & \textbf{Gew.} & \textbf{Score} & \textbf{Notiz} \\
\hline
1 & Differenzierung & 15\% & /10 & \\
\hline
2 & Taxonomie (AFB) & 10\% & /10 & \\
\hline
3 & Lernziel-Alignment & 10\% & /10 & \\
\hline
4 & Aktivierung & 10\% & /10 & \\
\hline
5 & Scaffolding & 10\% & /10 & \\
\hline
6 & Feedback-Qualität & 10\% & /10 & \\
\hline
7 & Sprachliche Klarheit & 10\% & /10 & \\
\hline
8 & Motivation \& Relevanz & 10\% & /10 & \\
\hline
9 & Inklusion (UDL) & 10\% & /10 & \\
\hline
10 & Praktikabilität & 5\% & /10 & \\
\hline
\multicolumn{3}{|r|}{\textbf{GESAMT:}} & \textbf{/10} & \\
\hline
\end{tabular}

\vspace{1em}
\textbf{Bewertungsschwellen:}
\begin{itemize}
    \item[$\Box$] \textcolor{successgreen}{$\geq$ 9.0: \textbf{FREIGABE} (Premium-Qualität)}
    \item[$\Box$] \textcolor{warningorange}{8.0--8.9: Iteration empfohlen}
    \item[$\Box$] 6.0--7.9: Überarbeitung erforderlich
    \item[$\Box$] $<$ 6.0: Grundlegende Neukonzeption
\end{itemize}
\end{scorebox}

\tierone{
\subsection{Score-Begründungen}

\textbf{1. Differenzierung (15\%)}
\begin{itemize}
    \item Qualitative Differenzierung (nicht nur mehr/weniger)?
    \item Passt zum Lerngruppen-Profil (\lerngruppenprofil)?
    \testdifferenziert{\item 3 Niveaus konsistent durchgehalten?}
    \testeinheitlich{\item Interne Differenzierung erkennbar?}
\end{itemize}

\textbf{2. Taxonomie/AFB (10\%)}
\begin{itemize}
    \item Verteilung ca. 25/50/25?
    \item Operatoren korrekt verwendet?
\end{itemize}

\textbf{3. Lernziel-Alignment (10\%)}
\begin{itemize}
    \item Passt zum Rahmenplan M-V?
    \item Kompetenzen adressiert?
\end{itemize}

\textbf{4. Aktivierung (10\%)}
\begin{itemize}
    \item Konstruktive/generative Aufgaben?
    \item Nicht nur reaktives Ankreuzen?
\end{itemize}

\textbf{5. Scaffolding (10\%)}
\begin{itemize}
    \item Gestufte Hilfen vorhanden?
    \item Lösung nicht verraten?
\end{itemize}

\textbf{6. Feedback-Qualität (10\%)}
\begin{itemize}
    \item Elaboriertes Feedback?
    \item Prozessbezogene Rückmeldung?
\end{itemize}

\textbf{7. Sprachliche Klarheit (10\%)}
\begin{itemize}
    \item Altersgerechte Sprache (Jg. \jahrgangsstufe)?
    \item Fachbegriffe erklärt?
\end{itemize}

\textbf{8. Motivation \& Relevanz (10\%)}
\begin{itemize}
    \item Alltagsbezug vorhanden?
    \item ARCS-Elemente (Attention, Relevance, Confidence, Satisfaction)?
\end{itemize}

\textbf{9. Inklusion/UDL (10\%)}
\begin{itemize}
    \item Multiple Repräsentationen (Text + Bild + ...)?
    \item Barrierefreiheit (LRS, Förderbedarf)?
\end{itemize}

\textbf{10. Praktikabilität (5\%)}
\begin{itemize}
    \item Zeitrahmen realistisch (45 Min.)?
    \item Material vorhanden/erstellbar?
\end{itemize}
}

\tieronetwo{
\subsection{Verbesserungsmaßnahmen}
Falls Score $<$ 9.0, hier die Top-3 Maßnahmen:

\begin{enumerate}
    \item \textbf{[Schwächste Dimension]:} Konkrete Maßnahme...
    \item \textbf{[Zweitschlechteste]:} Konkrete Maßnahme...
    \item \textbf{[Drittschlechteste]:} Konkrete Maßnahme...
\end{enumerate}
}

% ============================================================================
% ENDE
% ============================================================================
\end{document}
